\documentclass[10pt,a4paper]{article}

% Shared presentation layer for the standalone R0--R7 development documents.
% Method equations remain authoritative in the canonical idea sketch.
\usepackage[a4paper,margin=18mm,headheight=14pt]{geometry}
\usepackage{fontspec}
\usepackage{xeCJK}
\xeCJKsetup{CJKspace=true}
\setmainfont{DejaVu Serif}
\setsansfont{DejaVu Sans}
\setmonofont{DejaVu Sans Mono}
\setCJKmainfont{Noto Serif CJK KR}
\setCJKsansfont{Noto Sans CJK KR}
\setCJKmonofont{Noto Sans Mono CJK KR}

\usepackage{amsmath,amssymb}
\usepackage{booktabs,longtable,tabularx,array}
\usepackage{enumitem}
\usepackage{xcolor}
\usepackage{hyperref}
\usepackage{fancyhdr}
\usepackage{microtype}
\usepackage[most]{tcolorbox}

\definecolor{CoreBlue}{HTML}{1F4E79}
\definecolor{CoreLight}{HTML}{EAF3F8}
\definecolor{StopRed}{HTML}{9C0006}
\definecolor{StopLight}{HTML}{FCE8E6}
\definecolor{GateGreen}{HTML}{38761D}
\definecolor{GateLight}{HTML}{EAF4E2}
\definecolor{DecisionOrange}{HTML}{B45F06}
\definecolor{DecisionLight}{HTML}{FFF2CC}

\hypersetup{
  colorlinks=true,
  linkcolor=CoreBlue,
  urlcolor=CoreBlue
}

\setlength{\parindent}{0pt}
\setlength{\parskip}{0.3em}
\setlength{\emergencystretch}{3em}
\setlist[itemize]{leftmargin=1.45em,itemsep=0.1em,topsep=0.2em}
\setlist[enumerate]{leftmargin=1.75em,itemsep=0.14em,topsep=0.22em}
\renewcommand{\arraystretch}{1.15}
\newcolumntype{L}[1]{>{\raggedright\arraybackslash}p{#1}}
\newcommand{\code}[1]{{\small\nolinkurl{#1}}}
\newcommand{\must}[1]{\textcolor{CoreBlue}{\textbf{#1}}}
\newcommand{\haltitem}[1]{\textcolor{StopRed}{\textbf{#1}}}
\newcommand{\passitem}[1]{\textcolor{GateGreen}{\textbf{#1}}}
\newcommand{\decisionitem}[1]{\textcolor{DecisionOrange}{\textbf{#1}}}

\newenvironment{contractbox}[1]
  {\begin{tcolorbox}[colback=CoreLight,colframe=CoreBlue,title={#1},fonttitle=\bfseries,before upper=\raggedright]}
  {\end{tcolorbox}}
\newenvironment{decisionbox}[1]
  {\begin{tcolorbox}[colback=DecisionLight,colframe=DecisionOrange,title={#1},fonttitle=\bfseries,before upper=\raggedright]}
  {\end{tcolorbox}}
\newenvironment{stopbox}[1]
  {\begin{tcolorbox}[colback=StopLight,colframe=StopRed,title={#1},fonttitle=\bfseries,before upper=\raggedright]}
  {\end{tcolorbox}}


\pagestyle{fancy}
\fancyhf{}
\lhead{Wind3DGS 개발 문서}
\rhead{R5 Local Residual}
\cfoot{\thepage}

\title{R5 개발 계약\\Local Necessity and Always-on Residual}
\author{Wind3DGS}
\date{2026-08-22}

\begin{document}
\maketitle

\begin{contractbox}{문서 역할}
이 문서는 frozen Best Global이 설명하지 못한 response가 실제로 localized Local package를
필요로 하는지 검증하고, 조건부 선택 이전의 learned always-on Local을 확정하기 위한 실행 계약이다.
R5는 Local necessity와 Gate C까지만 소유하며 temporal selector/state machine은 R6가 소유한다.
Canonical \code{eq:rd-local-residual}과
\code{eq:rd-local-support-projection}을 이름으로 참조한다.
\end{contractbox}

\tableofcontents

\section{목적과 소유권}

R5의 목적은 ``Local을 추가하면 좋아진다''가 아니라 다음 질문을 순서대로 닫는 것이다.

\begin{enumerate}
  \item Frozen Best Global의 held-out residual이 teacher/mapping noise보다 크고 공간적으로 localized되는가?
  \item Larger Global보다 support-constrained oracle Local이 actual matched-p95에서 유리한가?
  Same active-DOF에서도 같은 경향이 유지되는가는 별도 mandatory secondary diagnostic으로 확인한다.
  \item Oracle ceiling이 확인된 뒤 learned always-on Local이 그 이득을 재현하는가?
  \item 같은 support/budget의 causal procedural flutter보다 learned Local이 response를 더 잘 맞추는가?
\end{enumerate}

R5는 residual atlas, Local support package, oracle/learned/procedural 비교와 Gate C를 소유한다.
R4의 Global checkpoint를 다시 학습하거나 rank/config를 바꾸지 않으며, R6 selector의 latency claim을 하지 않는다.

\section{진입 입력}

\begin{longtable}{L{0.29\linewidth}L{0.59\linewidth}}
\toprule
입력 & 요구 내용\\
\midrule
\endfirsthead
\toprule
입력 & 요구 내용\\
\midrule
\endhead
\code{BestGlobalCheckpoint} & R4 Gate B가 동결한 network/rank/config/checkpoint와 dependency hash\\
\code{GlobalResponsePackageSet} & Best Global model이 declared asset마다 compile한 immutable Global package와 evaluator hash\\
\code{R5GlobalHandoff} & frozen Global rollout, common-probe residual extraction API와 larger-Global baseline registry\\
\code{TeacherPackage} & converged teacher package collection, force/work와 mapping-noise floor\\
\code{RendererBranchContract} & R0에서 동결한 \(B_R\) 또는 local-affine normal/covariance branch identity\\
\code{BaseTransportReport} & R1/R2 base/Global mean--normal--covariance transport와 traction fixture hash\\
\code{SupportProposalConfig} & static-GS-only proposal의 frozen input, rule/network, slot/rank/budget와 hash\\
\code{MetricRegistry} & residual locality, target-band amplitude/phase/PSD, stability, dispersion와 latency metric\\
\code{ThresholdRegistry} & oracle admission, learned Local success와 Gate C Boolean threshold\\
\code{ProceduralRegistry} & sine/noise band, causal envelope/phase, seed set와 matched-budget 규칙\\
\bottomrule
\end{longtable}

Gate B를 통과하지 못했거나 Best Global identity가 하나로 고정되지 않았으면 R5를 시작하지 않는다.
Support proposal은 teacher residual이나 test error를 보지 않는 static-GS-only 입력으로 test-open 전에 동결한다.

\section{이미 동결된 계약}

\subsection{Frozen Global residual}

Local label은 teacher full motion이 아니라 frozen Global이 설명한 component를 뺀 complementary response다.
\[
  r_L^*=r_{\mathrm{teacher}}-r_G^* .
\]
\(r\)은 canonical common-probe의 time/frequency response이며 개별 eigenmode column이 아니다.
Global network, field, pole, rank, checkpoint와 evaluator는 Local training 내내 freeze한다.
Local gradient가 Global parameter로 흐르거나 joint fine-tuning으로 residual definition이 이동하는 경로는 금지한다.

\subsection{Serialized support의 단일 정의}

각 Local slot은 처음부터 \textbf{free이면서 physically valid인 Gaussian만 포함하는}
하나의 ordered hard support
\[
  \mathcal I_k=(i_{k,1},\ldots,i_{k,n_k})
\]
를 package에 저장한다. ``raw support를 저장한 뒤 attachment/invalid set을 런타임에서 다시 빼는''
두 번째 support 정의를 두지 않는다. Support order, proposal identity, membership, mass, coverage,
rank와 hash는 serialization round trip 뒤에도 동일해야 한다.

\subsection{Support-internal Global-null construction}

\(J_k\in\{0,1\}^{3N\times3n_k}\)를 support-to-full injection,
\(\widetilde C_k=J_k^\top\widetilde B_{L,k}\),
\(M_k=J_k^\top M_xJ_k\), \(A_k=B_G^\top M_xJ_k\)라 둔다.
Raw Local candidate는 support 내부에서만 Global-null projection과 whitening을 수행한다.
\[
  \begin{aligned}
  \Pi_k&=I-M_k^{-1}A_k^\top
  (A_kM_k^{-1}A_k^\top)^\dagger A_k,\\
  \overline C_k&=\Pi_k\widetilde C_k,\\
  C_k&=\overline C_k
  (\overline C_k^\top M_k\overline C_k)^{-1/2},\\
  B_{L,k}&=J_kC_k,\qquad
  Q_{L,k}=C_k^\top J_k^\top F.
  \end{aligned}
\]
Accepted slot은 construction상 다음을 동시에 만족한다.
\[
  B_{L,k}[i,:]=0\quad(i\notin\mathcal I_k),\qquad
  B_G^\top M_xB_{L,k}=0,\qquad
  B_{L,k}^\top M_xB_{L,k}=I.
\]
Support가 free+valid row만 가지므로 attachment/invalid row도 exact zero다.
Dense Global complement 뒤 hard mask, 또는 hard mask 뒤 dense projection으로 locality/직교성 중 하나를
깨는 구현은 허용하지 않는다. Load projection과 transport는 support-only sparse path를 사용한다.
Local \(\Omega_{L,k},Z_{L,k}\)는 이 support-internal Global-null projection과 mass whitening을 마친
\textbf{최종 coordinate}의 pole이다. Raw candidate coordinate의 pole을 그대로 재사용하지 않고,
V0는 accepted \(C_k\)를 고정한 common-probe residual-response fit으로 식별한다.
Operator transform/re-diagonalization은 별도 method revision 없이는 혼용하지 않는다.

\subsection{Local angular/normal branch의 소유권}

R5는 \code{RendererBranchContract}를 바꾸지 않지만 learned Local motion에 필요한 branch-specific payload는 소유한다.
\(B_R\) branch라면 \(B_{R,L,k}\)를 \(B_{L,k}\)와 같은 final generalized coordinate, ordered support와
pole에 연결해 예측하고 serialize한다. Support 밖, attachment와 invalid angular row는 exact zero이며
독립 angular pole/state를 만들지 않는다. Local-affine branch라면 \(B_{R,L,k}\)를 만들지 않고,
Local이 이동시킨 mean에서 upstream versioned rule로 current normal/covariance를 계산한다.
어느 branch든 always-on Local rollout의 current normal이 다음 frame traction에 들어가는 경로까지 R5 fixture로 검증한다.

\subsection{검증 순서와 비교 공정성}

R5 구현 순서는 아래와 같이 고정한다.

\begin{enumerate}
  \item \textbf{Oracle Local:} larger Global sweep과 비교해 localized residual에 Local capacity가 필요한지 확인한다.
  \item \textbf{Learned always-on Local:} frozen static-GS proposal support에서 oracle response ceiling이
  개선된 경우에만 frozen-Global residual을 학습한다.
  \item \textbf{Procedural baseline:} learned package를 동결한 뒤 동일 support를 사용해
  same active-DOF 비교와 actual matched-p95 비교를 각각 실행한다.
\end{enumerate}

Oracle이 이득을 내지 못하면 learned Local을 시작하지 않는다. Learned always-on Local이 성공하기 전에는
procedural baseline 결과로 Local network architecture를 test별 튜닝하지 않는다. 최종 Gate C에는
frozen Global, larger Global, oracle Local, learned always-on Local과 procedural baseline을 모두 포함한다.

\section{구현 체크리스트}

\subsection{Residual atlas}

\begin{itemize}
  \item[$\square$] Frozen Global identity/hash 검사 후 teacher와 동일 common-probe/time grid에서 rollout
  \item[$\square$] Object/time/space/frequency별 \(r_L^*\)와 mapping-noise floor 저장
  \item[$\square$] Traveling/local gust에서 residual spatial localization과 target-band energy 측정
  \item[$\square$] 여러 object와 independent GS variant에서 residual band/region 반복성 검사
  \item[$\square$] Teacher spatial/temporal convergence error와 mapping floor 이하 residual은 Local label에서 분리
  \item[$\square$] Predeclared larger-Global rank/capacity sweep을 matched-p95 조건으로 실행
\end{itemize}

\subsection{Support proposal과 construction}

\begin{itemize}
  \item[$\square$] Static GS, typed attachment/validity와 immutable metadata만 proposal input으로 사용
  \item[$\square$] Proposal output을 free+valid ordered \(\mathcal I_k\) 하나로 정규화하고 package에 serialize
  \item[$\square$] Slot별 support size/radius, area/mass coverage, row 수와 proposal confidence 저장
  \item[$\square$] Support 내부 \(M_k,A_k,\Pi_k\) projection과 mass whitening 구현
  \item[$\square$] Pseudoinverse/Gram floor, retained rank/condition과 reject/reduce event 저장
  \item[$\square$] Accepted \(C_k\) 최종 coordinate에서 Local pole/damping을 식별하고 raw pole 재사용을 거부
  \item[$\square$] \(B_R\) branch이면 same-coordinate/support \(B_{R,L,k}\)를 구성하고
  support 밖/attachment/invalid angular row를 exact zero로 유지
  \item[$\square$] Local-affine branch이면 faded되지 않은 always-on Local mean에서 frozen rule로 current normal/covariance 계산
  \item[$\square$] Support 밖, attachment 및 invalid row numerical zero 검사
  \item[$\square$] \(B_G^\top M_xB_{L,k}=0\)와 \(B_{L,k}^\top M_xB_{L,k}=I\) fixture 통과
  \item[$\square$] Serialize/deserialize 후 support order, sparse field와 hash 동일성 검사
\end{itemize}

\subsection{Oracle에서 learned always-on Local까지}

\begin{itemize}
  \item[$\square$] Same-support oracle field를 만든 뒤 최종 projected/whitened coordinate의 pole을 common-probe residual로 fit
  \item[$\square$] Oracle과 larger Global의 same active-DOF 비교를 실행
  \item[$\square$] 별도 actual matched-p95 비교를 실행하고 두 축의 결과를 섞지 않음
  \item[$\square$] Oracle admission threshold를 통과한 뒤에만 Local package head 학습 시작
  \item[$\square$] Frozen Global residual만 label로 사용하고 Global gradient/update가 0인지 감사
  \item[$\square$] Mode-column alignment loss 대신 common-probe time/FRF/rollout response loss 사용
  \item[$\square$] All-slot support-only load projection, oscillator update와 sparse transport 구현
  \item[$\square$] Always-on rollout stability, attachment/locality/Global-null과 independent-resplat consistency 검사
\end{itemize}

\subsection{Cross-slot 진단과 procedural 비교}

\begin{itemize}
  \item[$\square$] \(G_{k\ell}=B_{L,k}^\top M_xB_{L,\ell}\),
  \(\rho_{k\ell}=\lVert G_{k\ell}\rVert_2\), support Jaccard와 combined effective rank 저장
  \item[$\square$] Simultaneous activation의 duplicate response와 amplitude overshoot 측정
  \item[$\square$] Locality를 깨는 global cross-slot QR과 새 coupled solver를 V0 경로에서 금지
  \item[$\square$] Procedural sine/noise가 learned Local과 정확히 같은 ordered support proposal set 사용
  \item[$\square$] Actual p95를 맞춘 primary 비교에서 load/update/transport 비용을 분리 profile
  \item[$\square$] 별도 same active-DOF secondary 비교를 실행하고 matched-p95 결과와 한 축으로 합치지 않음
  \item[$\square$] Causal envelope/phase가 현재와 과거 \(W/Q\)만 사용하고 future 정보가 없는지 검사
  \item[$\square$] Predeclared noise seed를 모두 aggregate하고 best seed 선택 금지
  \item[$\square$] Asset/test-result별 band, amplitude, support 또는 seed tuning 금지
\end{itemize}

\section{Open Design Decisions}

\begin{decisionbox}{핵심 미결정: support proposal}
Local support를 어떤 static-GS-only 규칙 또는 network로 제안할지는 아직 설계가 필요하다.
Proposal은 teacher residual/test error를 볼 수 없으며, test-open 전 freeze되어 oracle, learned Local과
procedural baseline이 동일한 ordered support set을 사용해야 한다.
\end{decisionbox}

\begin{enumerate}
  \item \decisionitem{Proposal 형식:} patch query, geometry/normal 기반 deterministic rule 또는
  setup-network support head 중 V0 방식을 선택한다.
  \item \decisionitem{Support budget:} slot 수, support radius/size, slot당 rank, overlap allowance와
  total sparse DOF budget을 정한다.
  \item \decisionitem{Free+valid 판정 소비:} R0 physical-valid contract를 support proposal 전/후 어느 API에서
  강제하고 empty/insufficient support를 어떻게 reject할지 정한다.
  \item \decisionitem{Rank failure:} pseudoinverse/Gram floor 아래일 때 fixed-rank reject 또는
  deterministic rank reduction 중 하나를 정한다. 결과를 보고 support proposal을 반복 생성하지 않는다.
  \item \decisionitem{Cross-slot 중복:} 진단 threshold와 반복적인 visible overshoot가 있을 때 사용할
  deterministic merge/exclusion 규칙을 정한다. 새 coupled solver는 V0 범위 밖이다.
  \item \decisionitem{Oracle construction:} V0 residual-response fit의 objective, frequency band,
  initialization, positivity와 실패 threshold를 정한다.
  \item \decisionitem{Gate C 수치:} Local target band, actual matched-p95의 primary
  amplitude/phase/PSD metric과 allowance, stability/dispersion veto 및 Boolean 식을 dev에서 동결한다.
  Same active-DOF는 선택 가능한 Gate 축이 아니라 mandatory secondary diagnostic으로 고정한다.
\end{enumerate}

\section{산출물}

\begin{longtable}{L{0.31\linewidth}L{0.57\linewidth}}
\toprule
산출물 & 필수 내용\\
\midrule
\endfirsthead
\toprule
산출물 & 필수 내용\\
\midrule
\endhead
\code{GlobalResidualAtlas} & frozen Global hash, object/space/time/frequency residual와 noise-floor 분해\\
\code{SupportManifest} & free+valid ordered support, proposal/config, rank/coverage와 serialization hash\\
\code{OracleLocalReport} & same-support oracle fit, larger-Global 비교, admission Boolean과 failure reason\\
\code{AlwaysOnLocalCheckpoint} & frozen Global dependency, Local model weights/config와 schema hash; per-asset package는 포함하지 않음\\
\code{AlwaysOnLocalPackageSet} & declared asset별 ordered support, translational/branch-specific angular field, pole와 evaluator를 직렬화한 package set\\
\code{LocalConstructionReport} & locality, attachment/validity/angular zero, Global-null, Gram과 cross-slot diagnostics\\
\code{ProceduralRegistry} & exact support, sine/noise band/envelope/seed, causal rule와 budget hash\\
\code{R5ProfilerFixture} & all-slot load projection/update/transport를 분리 측정하는 frozen hardware/precision/warm-up protocol\\
\code{GateCReport} & frozen/larger/oracle/learned/procedural 비교, actual p95와 Boolean 판정\\
\bottomrule
\end{longtable}

R6 hand-off는 Gate C를 통과한 \code{AlwaysOnLocalCheckpoint}와 \code{AlwaysOnLocalPackageSet},
actual all-slot cost와 selector가 사용할 허용 metadata만 포함한다. Model checkpoint와 per-asset package는
별도 hash로 유지하며 R5는 selector score/state를 생성하지 않는다.

\section{필수 fixture와 metric}

\subsection{필수 fixture}

\begin{enumerate}
  \item \textbf{Frozen owner:} Global parameter/optimizer gradient와 checkpoint hash 불변성
  \item \textbf{Support round trip:} ordered free+valid support와 sparse field serialization
  \item \textbf{Construction:} support 밖/attachment/invalid zero, Global-null과 within-slot mass Gram
  \item \textbf{Final-coordinate pole:} projected/whitened \(C_k\)와 pole의 coordinate identity 및 raw-pole 거부
  \item \textbf{Local angular/normal:} selected branch의 same-coordinate/support/zero/serialization과
  always-on current-normal traction replay
  \item \textbf{Cross-slot:} overlap/Jaccard, spectral coupling, effective rank와 simultaneous overshoot
  \item \textbf{Oracle ceiling:} localized gust 및 target band에서 larger Global 대비 response/latency Pareto
  \item \textbf{Learned always-on:} oracle support에서 held-out amplitude/phase/PSD와 long rollout
  \item \textbf{Procedural parity:} same support에서 actual matched-p95 primary와 별도 same active-DOF
  secondary run, causal history 및 all-seed aggregate
  \item \textbf{Sparse system path:} projection/update/transport 각각의 p50/p95/p99와 memory
\end{enumerate}

\subsection{Gate C metric}

\begin{itemize}
  \item residual energy가 teacher spatial/temporal convergence와 mapping-noise floor를 넘는 정도
  \item spatial localization, target-band amplitude, phase, PSD와 common-probe velocity error
  \item strong/larger Global 대비 oracle 및 learned always-on Local의 actual matched-p95 primary improvement
  \item 별도 same active-DOF mandatory secondary diagnostic
  \item learned Local 대비 same-support/budget procedural sine/noise의 paired difference
  \item long-rollout stability, independent-resplat response dispersion와 duplicate overshoot
\end{itemize}

Primary와 threshold는 test-open 전 \code{ThresholdRegistry}에 동결한다. Visual improvement만으로
Gate C를 통과시키지 않으며, setup time/package memory/throughput도 mandatory secondary diagnostic으로 보고한다.

\section{종료 조건과 실패 경로}

\begin{contractbox}{R5 종료 조건}
Residual atlas가 noise floor 위의 localized failure를 입증하고, oracle Local이 larger Global보다
dev-frozen actual matched-p95 primary metric에서 이득을 보이며, learned always-on Local이 같은 support에서 그 이득을 재현하고
procedural baseline보다 frozen primary metric을 개선하면 Gate C를 통과한다.
이때만 immutable always-on package를 R6로 넘긴다.
\end{contractbox}

\begin{itemize}
  \item \passitem{Gate C pass:} always-on Local을 core에 유지하고 R6 conditional execution으로 진행한다.
  \item \haltitem{Residual 없음:} residual이 teacher/mapping floor 이하이거나 localized/repeatable하지 않으면
  Local 구현을 중단하고 Global-only 논문으로 축소한다.
  \item \haltitem{Oracle failure:} larger Global보다 matched-p95 이득이 없으면 learned Local을 학습하지 않는다.
  \item \haltitem{Learned failure:} oracle은 성공하지만 learned always-on이 재현하지 못하면
  Local network/label/support 설계를 진단하되 selector로 이동하지 않는다.
  \item \haltitem{Procedural parity failure:} learned Local이 same-support/budget procedural baseline보다
  primary metric을 개선하지 못하면 Gate C를 실패 처리하고 Global-only 경로로 종료한다.
  \item \haltitem{Construction failure:} locality, attachment/validity zero 또는 Global-null을 동시에 보존하지 못하면
  dense projection이나 새 solver로 우회하지 않고 support/rank proposal을 재설계한다.
  \item \haltitem{Gate C failure:} Local과 selector claim을 모두 중단하고 R4의 frozen Global-only 결과로 종료한다.
\end{itemize}

\end{document}
